% !TEX root = ../algebra_02.tex

\section{Матрицы линейных отображений и собственные значения}

\subsection{Матрица линейного отображения}

Пусть $V,W$ --- конечномерные линейные пространства над полем $K$,
$\dim V=n$, $\dim W=m$.

Пусть $E=(e_1,\ldots,e_n)$ --- базис $V$, а $F=(f_1,\ldots,f_m)$ --- базис $W$.

\begin{Definition}{Координатный столбец}{}
	Для $v\in V$ координатный столбец $[v]_E\in K^n$ определяется равенством
	$v=E[v]_E$, то есть
	\[
	v=e_1x_1+\cdots+e_nx_n
	\quad \Longleftrightarrow \quad
	[v]_E=
	\begin{pmatrix}
		x_1\\
		\vdots\\
		x_n
	\end{pmatrix}.
	\]
\end{Definition}

\begin{Definition}{Матрица линейного отображения}{}
	Пусть $\alpha\colon V\to W$ --- линейное отображение. Его матрицей в базисах $E,F$ называется матрица $[\alpha]_{E,F}\in M_{m\times n}(K)$, столбцы которой равны координатам образов базисных векторов:
	\[
	[\alpha]_{E,F}
	=
	\bigl(
	[\alpha e_1]_F,\ldots,[\alpha e_n]_F
	\bigr).
	\]
	Иначе говоря,
	\[
	\alpha E=F[\alpha]_{E,F}.
	\]
\end{Definition}

\begin{Statement}{Координаты образа вектора}{}
	Пусть $\alpha\colon V\to W$ --- линейное отображение и $A=[\alpha]_{E,F}$. Тогда для любого $v\in V$
	\[
	[\alpha v]_F=A[v]_E.
	\]
\end{Statement}

\begin{proof}
	Так как $v=E[v]_E$, то
	\[
	\alpha v=\alpha(E[v]_E)=(\alpha E)[v]_E
	=(FA)[v]_E=F(A[v]_E).
	\]
	Значит, $[\alpha v]_F=A[v]_E$.
\end{proof}

\subsection{Композиция линейных отображений}

Пусть заданы линейные отображения
\[
V \xrightarrow{\alpha} W \xrightarrow{\beta} U,
\]
а также базисы $E$ в $V$, $F$ в $W$ и $G$ в $U$.

\begin{Theorem}{Матрица композиции}{}
	Для композиции $\beta\alpha\colon V\to U$ выполнено
	\[
	[\beta\alpha]_{E,G}
	=
	[\beta]_{F,G}[\alpha]_{E,F}.
	\]
\end{Theorem}

\begin{proof}
	Пусть $[\alpha]_{E,F}=A$, $[\beta]_{F,G}=B$. Тогда
	$\alpha E=FA$ и $\beta F=GB$. Поэтому
	\[
	(\beta\alpha)E
	=
	\beta(\alpha E)
	=
	\beta(FA)
	=
	(\beta F)A
	=
	(GB)A
	=
	G(BA).
	\]
	Следовательно, $[\beta\alpha]_{E,G}=BA=[\beta]_{F,G}[\alpha]_{E,F}$.
\end{proof}

\subsection{Замена базисов}

\begin{Definition}{Матрица перехода}{}
	Пусть $E,E'$ --- два базиса одного пространства $V$.
	Матрицей перехода от $E'$ к $E$ называется матрица
	\[
	C_{E'\to E}:=[\operatorname{id}_V]_{E',E}.
	\]
	Она переводит координаты в базисе $E'$ в координаты в базисе $E$:
	$[v]_E=C_{E'\to E}[v]_{E'}$.
\end{Definition}

\begin{Theorem}{Формула замены базисов}{}
	Пусть $\alpha\colon V\to W$, $E,E'$ --- базисы $V$, а $F,F'$ --- базисы $W$.
	Тогда
	\[
	[\alpha]_{E',F'}
	=
	C_{F\to F'}[\alpha]_{E,F}C_{E'\to E}.
	\]
\end{Theorem}

\begin{proof}
	Используем формулу для композиции:
	\[
	\alpha
	=
	\operatorname{id}_W\circ \alpha\circ \operatorname{id}_V.
	\]
	Тогда
	\[
	[\alpha]_{E',F'}
	=
	[\operatorname{id}_W]_{F,F'}
	[\alpha]_{E,F}
	[\operatorname{id}_V]_{E',E}
	=
	C_{F\to F'}[\alpha]_{E,F}C_{E'\to E}.
	\]
\end{proof}

\begin{Corollary}{Замена базиса для оператора}{}
	Пусть $\alpha\in \End V$, $E,E'$ --- базисы $V$, $[\alpha]_E=A$ и $C=C_{E'\to E}$.
	Тогда
	\[
	[\alpha]_{E'}=C^{-1}AC.
	\]
\end{Corollary}

\begin{Definition}{Подобные матрицы}{}
	Матрицы $A,A'\in M_n(K)$ называются подобными, если существует $C\in GL_n(K)$ такое, что
	$A'=C^{-1}AC$.
\end{Definition}

\begin{Remark}{Смысл подобия}{}
	Подобные матрицы --- это матрицы одного и того же линейного оператора в разных базисах.
\end{Remark}

\subsection{Нормальный вид матрицы линейного отображения}

\begin{Theorem}{Нормальный вид линейного отображения}{}
	Пусть $\alpha\colon V\to W$ --- линейное отображение,
	$\dim V=n$, $\dim W=m$, $\rk\alpha=r$.
	Тогда существуют базисы $E$ в $V$ и $F$ в $W$ такие, что
	\[
	[\alpha]_{E,F}
	=
	\begin{pmatrix}
		I_r & 0\\
		0 & 0
	\end{pmatrix}.
	\]
\end{Theorem}

\begin{proof}
	Выберем базис $f_1,\ldots,f_r$ пространства $\Im\alpha$.
	Для каждого $i=1,\ldots,r$ выберем $e_i\in V$ так, что $\alpha e_i=f_i$.
	
	Пусть $e_{r+1},\ldots,e_n$ --- базис $\Ker\alpha$.
	Тогда $e_1,\ldots,e_n$ --- базис $V$.
	
	Действительно, если
	$a_1e_1+\cdots+a_re_r+b_{r+1}e_{r+1}+\cdots+b_ne_n=0$,
	то после применения $\alpha$ получаем
	$a_1f_1+\cdots+a_rf_r=0$.
	Значит, $a_1=\cdots=a_r=0$, а затем и все $b_i=0$.
	
	Дополним $f_1,\ldots,f_r$ до базиса $F=(f_1,\ldots,f_m)$ пространства $W$.
	В базисах $E=(e_1,\ldots,e_n)$ и $F$ имеем:
	$\alpha e_i=f_i$ при $i\le r$ и $\alpha e_i=0$ при $i>r$.
	Отсюда получается указанная матрица.
\end{proof}

\begin{Corollary}{Ранг линейного отображения и ранг матрицы}{}
	Для любых базисов $E,F$ выполнено
	\[
	\rk \alpha=\rk [\alpha]_{E,F}.
	\]
\end{Corollary}

\subsection{Пространство линейных отображений}

\begin{Definition}{Пространство $\Hom(V,W)$}{}
	Множество всех линейных отображений из $V$ в $W$ обозначается через $\Hom(V,W)$.
	
	Оно является линейным пространством относительно операций
	\[
	(\alpha+\beta)(v)=\alpha v+\beta v,
	\qquad
	(\lambda\alpha)(v)=\lambda\alpha v.
	\]
\end{Definition}

\begin{Theorem}{Матрицы как координаты линейных отображений}{}
	Пусть $E$ --- базис $V$, $F$ --- базис $W$.
	Отображение
	\[
	\Phi\colon \Hom(V,W)\to M_{m\times n}(K),
	\qquad
	\Phi(\alpha)=[\alpha]_{E,F},
	\]
	является изоморфизмом линейных пространств.
\end{Theorem}

\begin{proof}
	Линейность следует из равенств
	\[
	[\alpha+\beta]_{E,F}=[\alpha]_{E,F}+[\beta]_{E,F},
	\qquad
	[\lambda\alpha]_{E,F}=\lambda[\alpha]_{E,F}.
	\]
	
	Покажем, что $\Phi$ биективно.
	Матрица $[\alpha]_{E,F}$ полностью определяется векторами
	$\alpha e_1,\ldots,\alpha e_n$.
	
	Обратно, если дана произвольная матрица $A=(a_1,\ldots,a_n)$,
	где $a_i\in K^m$ --- её столбцы, то можно единственным образом задать линейное отображение $\alpha$ формулами
	$\alpha e_i=Fa_i$.
	Тогда $[\alpha]_{E,F}=A$.
\end{proof}

\begin{Corollary}{Размерность пространства линейных отображений}{}
	Если $\dim V=n$, $\dim W=m$, то
	\[
	\dim \Hom(V,W)=mn.
	\]
\end{Corollary}

\subsection{Линейные операторы и матричная алгебра}

\begin{Definition}{Линейный оператор}{}
	Линейным оператором на $V$ называется линейное отображение из $V$ в себя.
	Множество всех таких операторов обозначается
	\[
	\End V:=\Hom(V,V).
	\]
\end{Definition}

\begin{Proposition}{Алгебра операторов и матриц}{}
	Пусть $E$ --- базис $V$, $\dim V=n$.
	Отображение
	\[
	\xi_E\colon \End V\to M_n(K),
	\qquad
	\xi_E(\alpha)=[\alpha]_E,
	\]
	является изоморфизмом алгебр:
	\[
	[\alpha+\beta]_E=[\alpha]_E+[\beta]_E,
	\qquad
	[\lambda\alpha]_E=\lambda[\alpha]_E,
	\qquad
	[\beta\alpha]_E=[\beta]_E[\alpha]_E.
	\]
\end{Proposition}

\begin{Remark}{Единица в $\End V$}{}
	В алгебре $\End V$ единицей является тождественный оператор $\operatorname{id}_V$.
	Ему соответствует единичная матрица $I_n$.
\end{Remark}

\subsection{Инвариантные подпространства}

\begin{Definition}{Инвариантное подпространство}{}
	Пусть $\alpha\in \End V$ и $W\le V$.
	Подпространство $W$ называется инвариантным относительно $\alpha$, если
	\[
	\alpha W\subset W.
	\]
\end{Definition}

\begin{Theorem}{Блочный вид при инвариантном подпространстве}{}
	Пусть $W\le V$, $\dim W=m$, и
	$E=(e_1,\ldots,e_n)$ --- базис $V$ такой, что
	$e_1,\ldots,e_m$ --- базис $W$.
	
	Тогда $W$ инвариантно относительно $\alpha$ тогда и только тогда, когда
	\[
	[\alpha]_E=
	\begin{pmatrix}
		B & *\\
		0 & C
	\end{pmatrix},
	\]
	где $B\in M_m(K)$.
\end{Theorem}

\begin{proof}
	Если $W$ инвариантно, то для $j=1,\ldots,m$ имеем $\alpha e_j\in W$.
	Значит, в координатах по базису $E$ у вектора $\alpha e_j$ последние $n-m$ координат равны нулю.
	Именно поэтому в первых $m$ столбцах матрицы снизу стоит нулевой блок.
	
	Обратно, если матрица имеет такой блочный вид, то для $j=1,\ldots,m$ вектор $\alpha e_j$ лежит в $\langle e_1,\ldots,e_m\rangle=W$.
	Следовательно, $\alpha W\subset W$.
\end{proof}

\begin{Corollary}{Инвариантность прямых слагаемых}{}
	Пусть $V=W_1\oplus W_2$,
	$e_1,\ldots,e_m$ --- базис $W_1$, а $e_{m+1},\ldots,e_n$ --- базис $W_2$.
	Для $E=(e_1,\ldots,e_n)$ следующие условия эквивалентны:
	
	\begin{enumerate}
		\item $W_1$ и $W_2$ инвариантны относительно $\alpha$;
		\item матрица оператора имеет блочно-диагональный вид
		\[
		[\alpha]_E=
		\begin{pmatrix}
			B & 0\\
			0 & C
		\end{pmatrix}.
		\]
	\end{enumerate}
\end{Corollary}

\begin{Definition}{Индуцированный оператор на факторпространстве}{}
	Пусть $W$ инвариантно относительно $\alpha\in\End V$.
	Тогда на факторпространстве $V/W$ корректно задаётся оператор
	\[
	\overline{\alpha}\colon V/W\to V/W,
	\qquad
	\overline{\alpha}(v+W)=\alpha v+W.
	\]
\end{Definition}

\begin{proof}
	Проверим корректность. Если $v+W=v'+W$, то $v-v'\in W$.
	Так как $W$ инвариантно, $\alpha(v-v')\in W$.
	Значит, $\alpha v+W=\alpha v'+W$.
\end{proof}

\begin{Remark}{Связь блоков с ограничением и фактороператором}{}
	В условиях теоремы о блочном виде блок $B$ является матрицей ограничения
	$\alpha|_W\colon W\to W$ в базисе $e_1,\ldots,e_m$.
	
	Если $\overline e_{m+1},\ldots,\overline e_n$ --- базис $V/W$, где
	$\overline e_i=e_i+W$, то блок $C$ является матрицей индуцированного оператора
	$\overline{\alpha}$.
\end{Remark}

